\subsection{Strong integrability for the coherent branch}\label{lp:coherent}

The other analytic branch also admits an integrability improvement. Unlike the
frequency argument, this follows from the positive affine approximations.

\begin{proposition}\label{lp:coherent-prop}
Let \(\mu,\nu\) be probabilities on positively separated compact planar sets.
Suppose \(\mu\) is \(s\)-Frostman, its occupied-square count is
\(O(r^{-u})\), and \(\nu\) has a Frostman exponent greater than one, where
\(1<s\le2\) and \(s\le u<2s-1\).
Let \(q>1\) be the averaged angular exponent in the coherent center
construction. Put
\[
 \eta=\frac{q-1}{2q-1},\qquad \delta=2s-1-u.
\]
The joint law of \((y,|x-y|)\) under \(\mu\times\nu\) has a density
in \(L^p(d\nu\,dt)\) for every
\begin{equation}\label{lp:coherent-range}
 1<p<\left(1-\eta\frac{\delta}{s}\right)^{-1}.
\end{equation}
\end{proposition}
\begin{proof}
Set \(v=2-1/q\) and \(q'=q/(q-1)\), so
\(v-1=1/q'\) and \((v-1)/v=\eta\). Retain the chosen centers and
positive densities \(f_n\) from Section~\ref{sec:coherenttree}. For each
occupied square \(Q\) at scale \(r=2^{-n}\), let \(a_Q(y,t)\) be
its conditional affine density, and let \(h_Q(\theta,t)\) be the
orthogonal projection density before translation in the scalar variable.
The source weights and angular densities satisfy
\[
 p_Q=\mu(Q),\quad \mu_Q=p_Q^{-1}\mu|_Q,\quad
 F_Q=\|\rho_{x_Q}\|_q^q,\qquad \sum_Qp_QF_Q\le2B.
\]
Probability interpolation and then angular H\"older give
\[
 \int |h_Q(\theta,t)|^v\,dt
 \le\|h_Q(\theta,\cdot)\|_2^{2(v-1)},\qquad
 \|a_Q\|_v^v\le C F_Q^{1/q}I_1(\mu_Q)^{1/q'}.
\]
For the second inequality, use the angular pin law
\(\rho_{x_Q}(\theta)\,d\theta\), the equality \((v-1)q'=1\),
and the projection-energy identity. The pin-dependent translation preserves
each scalar norm. The source Frostman bound inside \(Q\) gives
\(I_1(\mu_Q)\le Cr^{s-1}/p_Q\).
Since \(f_n=\sum_Qp_Qa_Q\), convexity and H\"older over the
at most \(Cr^{-u}\) occupied squares yield
\begin{align}
 \|f_n\|_v^v
 &\le C r^{(s-1)/q'}\sum_Q(p_QF_Q)^{1/q}
 \le C B^{1/q}r^{(s-1-u)/q'},\label{lp:coherent-growth}\\
 \|f_n\|_v&\le C r^{-A\eta},\qquad A=u+1-s>0.\nonumber
\end{align}
No lower bound on an individual square mass was used.

Choose \(0<\gamma<\min\{1/2,(s-1)/2\}\) with
\(\tau=s-u+2\gamma>0\). The optimized coherent comparison and
the covering estimate in Proposition~\ref{prop:coherentcovering} give
\[
 \|f_{n+1}-f_n\|_1\le C r^{\tau\eta},\qquad
 \|f_{n+1}-f_n\|_v\le C r^{-A\eta}.
\]
For \(1<p<v\), set
\(\vartheta=(v-p)/(p(v-1))\), so
\(1/p=\vartheta+(1-\vartheta)/v\). H\"older interpolation gives
\[
 \|f_{n+1}-f_n\|_p
 \le C r^{\eta[\tau\vartheta-A(1-\vartheta)]}.
\]
The exponent is positive exactly when
\[
 p<\frac{v(\tau+A)}{\tau+vA}
 =\frac{v(1+2\gamma)}{s-u+2\gamma+v(u+1-s)}.
\]
Letting \(\gamma\) approach \((s-1)/2\) from below makes this upper
bound tend to
\[
 \frac{vs}{\delta+v(s-\delta)}
 =\left(1-\eta\frac\delta s\right)^{-1}.
\]
Thus every \(p\) in \eqref{lp:coherent-range} admits a fixed legal
\(\gamma\) giving summable dyadic differences in \(L^p\).
The common pin-distance region has finite measure, and an initial
approximation has finite \(v\)-norm by \eqref{lp:coherent-growth}.
Consequently \(f_n\) converges in \(L^p\) and \(L^1\).
The uniform affine approximation to distance identifies this limit with
the raw positive joint distance law, as in Theorem~\ref{thm:coherentcriterion}.
\end{proof}

\begin{theorem}[A stronger conclusion under the same dimension curve]
\label{lp:dimension}
Under the hypotheses of Theorem~\ref{target:hardgap}, there are separated
compact source and pin subsets of \(E\), carrying probabilities \(\mu,\nu\),
and an exponent \(p>1\), such that the raw joint distance law has a
density \(f\in L^p(d\nu\,dt)\). In particular
\((d_y)_*\mu\) has an \(L^p\) density for \(\nu\)-almost every pin.
\end{theorem}
\begin{proof}
In the coherent case, use the same strict Frostman and covering exponent
choices and separated compact restrictions as in
Proposition~\ref{prop:coherentborel}, followed by
Proposition~\ref{lp:coherent-prop}. In the other case, retain the
strict profile-cost margin and compact restrictions in the proof of
Theorem~\ref{target:hardgap}, and apply the preceding subquadratic
frequency argument. Both give a strong exponent greater than one for
the actual joint law. Fubini and its probability disintegration give
the final assertion. No uniform exponent as the strict margin tends
to zero is claimed.
\end{proof}

This also gives a quantitative length statement for a positive proportion of
the selected pins. If \(M=\int\!\int f(y,t)^p\,dt\,d\nu(y)\), then
pins in a set of \(\nu\)-measure at least \(1/2\) satisfy
\[
 |\Delta_y(E)|\ge(2M)^{-1/(p-1)}.
\]
Indeed Markov bounds \(\int f(y,t)^pdt\) by \(2M\) on such a pin
set, and H\"older applied to a density of mass one proves the lower
bound. This improves integrability, not the dimensional cutoff or its
strict endpoints.
